\begin{statement}{F.9}{Proved}{The genuine integral class-two action for the colliding-cusp curve}
Let $K_0=\mathbb Q((t))$ and
$U_t=\mathbb P^1_{K_0}\setminus\{0,t,1,\infty\}$.
Use the coefficient splitting of $G_{K_0}$ which fixes
compatible roots of $t$.  There is a based splitting
of the arithmetic fundamental-group extension,
constructed from the node chart below, for which
the free class-two pro-$\ell$ group has a basis
\[
                 x=b,\qquad y=d=ab,\qquad z=c
                                                               \tag{F.9.1}
\]
with the following genuine continuous action:
\[
 \begin{array}{c|ccc}
       &x&y&z\\ \hline
 \tau&yxy^{-1}&y&z\\
 g\in G_{\mathbb Q}&x^{\chi_\ell(g)}&y^{\chi_\ell(g)}
                                         &z^{\chi_\ell(g)}.
 \end{array}                                                 \tag{F.9.2}
\]
The inverse free-basis change is
$a=yx^{-1}$, $b=x$, $c=z$.
These formulas hold integrally for every prime,
including two.  In the original generators, putting
$c_{ab}=[a,b]$ and $\lambda=\chi_\ell(g)$, the constant
action is
\[
 g(a)=a^\lambda c_{ab}^{-\binom\lambda2},\qquad
 g(b)=b^\lambda,\qquad g(c)=c^\lambda .                    \tag{F.9.3}
\]
The inertia action in that basis is exactly (F.7)
of FLAGGED\_149 F17.3.  The comparison changes the
free basis and specifies a splitting; it preserves
the arithmetic extension over $G_{K_0}$.
\end{statement}
\begin{proof}
Blow up the point $x=t=0$ of the marked projective
line model.  Near its single node use coordinates
$X=x$ and $V=t/x$, so $XV=t$.
The outer component is the tripod in coordinate
$X$, with node at zero and marked points at one
and infinity.  The bubble is the tripod in
coordinate $V$, with node at zero, the cusp $t$
at one, and the cusp $0$ at infinity.
The branch parameters $X,V$ have product exactly
$t$; no constant factor has been inserted.

The punctured normal-crossing node chart has
arithmetic tame group
\[
 \bigl(\widehat{\mathbb Z}(1)^2\bigr)\rtimes G_{\mathbb Q}
 \longrightarrow
 \widehat{\mathbb Z}(1)\rtimes G_{\mathbb Q}=G_{K_0},
 \quad ((r,s),g)\longmapsto(r+s,g).                         \tag{F.9.4}
\]
This is seen on the compatible covers extracting
roots of $X,V$; the base cover extracts roots of
$XV$.  Choose their products to be the specified
roots of $t$.  Coefficient automorphisms fix these
formal roots and act on their root-of-unity
multipliers by $\chi$.  The map in (F.9.4) then has
the continuous group section
$(r,g)\mapsto((r,0),g)$.
The homomorphism from the node chart's arithmetic
group to that of $U_t$, induced by the punctured
formal-neighborhood map, sends this group section
to an actual section over $G_{K_0}$.
The kernel direction $(r,-r)$ is the vanishing
loop $d$; its primitive abelianization proves
injectivity on this direction.  Choose its
orientation and the lift of $\tau$ so that the
associated Dehn twist conjugates the bubble
group by $d$ and fixes the outer group.
This is the annulus cut-and-glue calculation
which gives (F.6) in the attached return.

We specify why the component comparison respects
arithmetic.  Logarithmic specialization for this
proper marked semistable model is Theorem 3.4.6
of Stix's thesis \cite{BC_Stix2002}; the residue
characteristic is zero, so its prime-to-residue-
characteristic assertion applies to every
pro-$\ell$ quotient.  The arithmetic van Kampen
presentation is its equation (6.2.5), with
component comparison Lemma 6.2.8 and injective
edge maps Corollary 6.2.9.  The dual graph here
is a single edge, so the based geometric
amalgam is
\[
 \widehat F(d,c)*_{\langle d\rangle}
                     \widehat F(d^{-1},b).                 \tag{F.9.5}
\]
Its arithmetic amalgam uses the same coefficient
section and the same node root fiber functor
on both sides of (F.9.4).  Thus the continuous
maps from the two component groups, and their
common inertia inclusion, are $G_{\mathbb Q}$-
equivariant with these actual sections.
The common generator is primitive in both
vertex groups; eliminating its duplicate in
the amalgam gives the free basis $b,d,c$.
The same elimination on finite $\ell$-group
quotients proves the pro-$\ell$ comparison.

On each component use the standard path from
the unit tangent at zero towards the point
one.  For a tripod with these data the based
arithmetic action sends its first loop to its
$\lambda$-th power and its second to a conjugate
of its $\lambda$-th power by an element of the
commutator subgroup.  The latter assertion
is the normalized action in Wickelgren,
\S2.1, equation for the $G_k$-action
\cite{BC_Wickelgren2012}; its abelian path
cocycle is zero because the two leading
units are one.  In the class-two quotient
this conjugator is central.  The outer pair
$(d,c)$ and bubble pair $(d^{-1},b)$ therefore
both have the pure-power action in (F.9.2).
Compatibility in (F.9.4) proves that they
give the same based action on (F.9.5).
The Dehn twist gives the first row of
(F.9.2).

For clarity about the original peripheral
generators, the standard bubble loop at
infinity is $b^{-1}d$.  Conjugating its path
by $d$ gives the inertia generator
$a=db^{-1}$, so $ab=d$ exactly.  This is
the indicated Nielsen basis change, with
the displayed inverse; it is a continuous
automorphism of each free pro-$\ell$ group.
Both $a$ and $b$ are conjugated by $d$ under
the twist, proving (F.7).
Finally the class-two collection identity
\[
          (ab)^\lambda=a^\lambda b^\lambda
                          [a,b]^{-\binom\lambda2}
\]
gives (F.9.3).  The polynomial
$\binom\lambda2=\lambda(\lambda-1)/2$
is integral for $\lambda\in\mathbb Z_2$
and extends continuously from integers;
thus no division by a nonunit is being
performed in the group law.  This retains
the sign information which would be lost
by declaring the original pair $a,d$
simultaneously diagonal without a tangent
comparison at the prime two.
\end{proof}

